\begin{abstract}
O retificador de Vienna é uma das topologias mais utilizadas em estágios de 
conversão CA-CC, em especial em carregadores de veículos elétricos (VEs) de média
e alta potência. Nesse sentido, este trabalho revisa topologias de conversores com correção
de fator de potência (PFC) aplicadas a VEs e detalha a topologia do retificador de
Vienna trifásico a quatro fios (ponto de neutro aterrado) com fator de potência variável.
Tendo isso em mente, é apresentado um modelo por fase e uma estrutura de controle para o 
retificador citado anteriormente, na qual são desenvolvidos controladores proporcional-ressonantes (PR)
de corrente para cada fase, um controlador proporcional-integral (PI) para o controle da tensão do barramento CC e
um controlador PI para o controle do equilíbrio do ponto médio e geração de referência com 
fator de potência (FP) variável. Por fim, são apresentadas simulações de um conversor de \SI{21}{kW},
que demonstram a eficácia do modelo e da estratégia de controle propostos, obtendo um THD inferior a \SI{2}{\percent} após o controle.
\end{abstract}

\begin{IEEEkeywords}
retificador de Vienna, fator de potência variável, controle ressonante, carregadores de veículos elétricos.
\end{IEEEkeywords}